% hip-compute-settlement.tex
% Hanzo Improvement Proposal — Compute Settlement: the ledger for paid,
% attested compute on a Hanzo sovereign L1 over Lux.
% Number assigned when this lands in HIPs/.
\documentclass[11pt,a4paper]{article}

\usepackage[utf8]{inputenc}
\usepackage[T1]{fontenc}
\usepackage[margin=1in]{geometry}
\usepackage{amsmath,amssymb,amsthm}
\usepackage{booktabs}
\usepackage{longtable}
\usepackage{array}
\usepackage{enumitem}
\usepackage{xcolor}
\usepackage{microtype}
\usepackage[colorlinks=true,linkcolor=black,citecolor=black,urlcolor=blue]{hyperref}

\newtheorem{definition}{Definition}[section]
% `rule' is a LaTeX primitive, so the environment carrying normative Rules is
% named `norm'. The label prefix and the printed head both stay "Rule".
\newtheorem{norm}[definition]{Rule}
\newtheorem{proposition}[definition]{Proposition}
\newtheorem{remark}[definition]{Remark}

% Long identifiers are the norm here, so let the typewriter font break.
\newcommand{\code}[1]{\texttt{\hyphenchar\font=`\-\relax #1}}
\sloppy
\emergencystretch=3em

\newcommand{\Adm}{\mathrm{Adm}}
\newcommand{\bond}{\mathrm{bond}}
\newcommand{\esc}{\mathrm{escrow}}
\newcommand{\res}{\mathrm{reserved}}
\newcommand{\pay}{\mathrm{pay}}
\newcommand{\emit}{\mathrm{emit}}
\newcommand{\open}{\mathrm{open}}

\setlist[itemize]{itemsep=2pt,topsep=4pt}
\setlist[enumerate]{itemsep=2pt,topsep=4pt}
\setlist[description]{itemsep=3pt,topsep=4pt}

\title{\textbf{Compute Settlement}\\[4pt]
\large Credits in, attested work out, on a Hanzo sovereign L1 over Lux\\[6pt]
\large A Hanzo Improvement Proposal}

\author{Hanzo AI}
\date{}

\begin{document}
\maketitle

\begin{center}
\begin{tabular}{@{}ll@{}}
\toprule
Number      & assigned on merge into \code{HIPs/} \\
Type        & Standards Track \\
Category    & Core \\
Status      & Draft \\
Requires    & HIP-0024 (Hanzo sovereign L1), HIP-0095 (QoS challenge),\\
            & HIP-0096 (compute rewards), HIP-0106 (plugin contract, usage\\
            & warehouse), HIP-1313 (\code{/v1/usage}), HIP-1325 (\code{/v1/node});\\
            & LP-018 (network security modes), LP-77 (P-Chain warp payloads) \\
\bottomrule
\end{tabular}
\end{center}

\begin{abstract}
\noindent
A compute marketplace has two money flows and one question between them. Buyers
pay for use; operators are paid for work; and somewhere the chain must decide
whether a particular operator did a particular piece of work well enough to be
paid for it. This proposal specifies that decision and the ledger around it, as
a Hanzo sovereign L1 on Lux. It designs no consensus, no signature scheme and no
attestation format: Lux supplies all three, and the parts of them this chain
uses are cited to the files that implement them.

The design rests on three separations. \emph{Who may claim} is answered by a
bond. \emph{What ran} is answered by a measurement and a piece of signed
evidence about it. \emph{What it is worth} is answered by a price a
counterparty already agreed to and escrowed. Braiding any two of these produces
the failure modes that dominate compute-market designs, and we keep them apart
in the state, in the transaction surface, and in the reward formula.

Two consequences are worth stating up front. First, emission is weighted by
\emph{paid} work rather than by self-reported capacity, and is rate-capped
rather than pooled; under those two conditions collusion between a buyer and an
operator is strictly unprofitable whenever the emission rate is below the
protocol fee rate, which we prove. Second, an invalid attestation is a
rejection and not a slash; only an accepted obligation that produces no admitted
claim is slashable. One rule replaces two, and the bond is sized so that every
slash it authorises is collectible.

We report what the current implementation does and does not do. The Lux
attestation verifier is real and fails closed, but the compute-market
precompile's own \code{verifyCompute} accepts any caller-computable hash, and
we give the fix. On the reference fleet, measured, neither host can produce
hardware-rooted evidence today, so the chain must be able to price two trust
bases rather than one.
\end{abstract}

\tableofcontents

\section{What the chain must decide}

Between a buyer who pays and an operator who is paid sits a question neither
can settle alone: what actually ran, and is it worth what was agreed? A
settlement chain exists to answer that question and to move money on the
answer. It is worth being exact about which parts of the question a chain can
answer at all.

Three properties are routinely conflated:

\begin{description}[leftmargin=1.4em]
  \item[Integrity of the inputs.] The files fed to the machine were the files
    they were said to be.
  \item[Fidelity of execution.] Those files, and only those, determined what
    ran.
  \item[Quality of the result.] What ran produced a correct answer.
\end{description}

They have different mechanisms. A digest over inputs addresses the first. A
hardware root of trust addresses the second. Neither addresses the third; that
is the domain of verification schemes over the computation itself and of
challenge protocols, and this proposal cites rather than restates them
(HIP-0095, LP-313).

This proposal specifies a chain that decides the first, prices the second
according to which root of trust is offered, and declines to assert the third.
Declining is not a gap to be papered over: a settlement rule that pretends to
witness result quality prices a guarantee it cannot deliver, and LP-324 is an
audit of what happens when a market does that.

\subsection{Scope}

In scope: the credit ledger, operator bonds and slashing, the admission rule
for a claim of served compute, the settlement arithmetic, and the exact
boundary between chain state and fleet state.

Out of scope, and specified elsewhere: consensus (Lux Quasar), the attestation
format and its verifier (Lux \code{precompile/ai}), inference verification
(LP-313), the challenge protocol over disputed quality (HIP-0095), price
discovery, and routing.

\subsection{Conventions}

The key words MUST, MUST NOT, SHOULD, SHOULD NOT and MAY are to be interpreted
as in RFC 2119. Statements are labelled by epistemic status throughout:
\emph{implemented} means the behaviour exists in a cited file; \emph{measured}
means we ran it and report the result; \emph{proposed} means this document is
asking for it; \emph{assumed} means a security or economic premise that is not
checked by the protocol.

\section{What exists today}

This section reports the substrate with locations, so every later claim can be
checked against code rather than against prose. Everything in it is
\emph{implemented} unless marked otherwise.

\subsection{The sovereign-L1 model}
\label{sec:sovereign}

Lux's P-Chain models a network's security along two orthogonal axes
(\code{node/vms/platformvm/security/mode.go}):

\begin{itemize}
  \item \code{RestakeParent} --- the network leans on its parent's validator set.
  \item an own set, described by \code{Admission} $\in$
    \{\code{NoOwnSet}, \code{Open}, \code{Gated}\} and \code{Manager} $\in$
    \{\code{PChain}, \code{Contract}\}, with a \code{Threshold} that is the
    minimum stake under \code{Open} admission.
\end{itemize}

\code{Mode.Sovereign()} is \code{Admission != NoOwnSet}. \code{Mode.Valid()}
enforces the cross-axis invariant that a network must either restake its parent
or run its own set; a network that does neither is refused with
\code{ErrNoSecurity}.

A network is born by \code{CreateNetworkTx}
(\code{node/vms/platformvm/txs/create\_network\_tx.go}), which is the sole
network constructor at any level of the hierarchy: \code{Parent = Primary}
yields an L1, \code{Parent =} an L1 yields an L2, and the wire is byte-identical
at every level. It carries the \code{security.Mode} inline (offsets
\code{offCN\_RestakeParent}, \code{offCN\_Admission}, \code{offCN\_Manager},
\code{offCN\_Threshold}), the genesis validator list, and
\code{ManagerChainID}/\code{ManagerAddress} naming the chain and address hosting
a \code{Contract}-governed set's staking contract.

Chains are not created by that transaction. \code{CreateChainTx} is the sole
chain constructor, and an L1 spawn is a block containing one
\code{CreateNetworkTx} and $N$ \code{CreateChainTx}, atomic at the block.

\begin{remark}
The brief for this work described L1 birth as
\code{CreateChainTx} + \code{AddValidatorTx} + \code{ConvertNetworkToL1Tx}.
That is the older shape. In the current tree the legacy
\code{ConvertNetworkToL1Tx} and \code{CreateSovereignL1Tx} have been folded into
\code{CreateNetworkTx}
(\code{node/vms/platformvm/network/zap\_native\_admission.go}), so a network may
be \emph{born} sovereign in one transaction. \code{ConvertNetworkTx}
(\code{convert\_network\_tx.go}) survives as the distinct endomorphism
$\text{Network} \to \text{Network}$: it promotes an existing network from
inherited to sovereign security and re-anchors its parent. Hanzo needs the
constructor, not the promotion, because it has no restaked history to promote.
\end{remark}

The L1 validator lifecycle is four transactions:
\code{RegisterL1ValidatorTx} (a funding balance, a BLS proof of possession, and
a signed Warp \code{AddressedCall} payload),
\code{IncreaseL1ValidatorBalanceTx} (top up the continuous-fee balance),
\code{DisableL1ValidatorTx} (authorised by credential over a
\code{ValidationID}), and the weight updates carried by the
\code{L1ValidatorWeight} Warp payload
(\code{node/vms/platformvm/warp/message/l1\_validator\_weight.go}, 49 bytes:
kind, \code{ValidationID}, nonce, weight). That payload is bidirectional --- the
P-Chain accepts it as a command and emits it as a report --- and a nonce of
$2^{64}-1$ with weight $0$ is reserved for removal. These payloads are specified
in LP-77.

\subsection{Consensus}
\label{sec:consensus}

Quasar is the Lux post-quantum finality binding
(\code{consensus/protocol/quasar}). It runs up to three independent signing
paths in parallel: BLS12-381 threshold signatures (ECDL), Corona Ring-LWE
two-round threshold signatures (Module-LWE), and ML-DSA-65 identity signatures
(FIPS 204; Module-LWE and Module-SIS). Each layer is independently toggleable,
so \{BLS\}, \{BLS, ML-DSA\}, \{BLS, Corona\} and all three are configurations of
one engine. Inter-node transport is ZAP with optional PQ-TLS 1.3.

This proposal makes no consensus choice. It requires only two properties that
Quasar already provides and that the settlement rules below depend on:
deterministic finality of a block height, and a validator set whose weight the
P-Chain can be told to change.

\subsection{The attestation verifier}
\label{sec:verifier}

There is one real TEE verifier in the AI precompile block:
\code{ai.VerifyTEE} $\to$ \code{verifyAttestation}
(\code{precompile/ai/tee\_attest.go}). Its receipt format is a 48-byte report
header followed by a concatenated DER certificate chain:

\[
  \underbrace{\text{deviceID}}_{32} \;\|\;
  \underbrace{\text{timestamp}}_{8,\ \text{BE}} \;\|\;
  \underbrace{\text{nonce}}_{8} \;\|\;
  \underbrace{\text{chain}_{\mathrm{DER}}}_{\text{leaf, intermediates}}
\]

and the signature is over the 48-byte header by the leaf key. Verification
checks that the chain terminates at an embedded vendor root, valid \emph{at the
report's own embedded timestamp} rather than the wall clock, and that the
signature binds the header. It is deterministic, so every validator reaches the
same verdict; it fails closed, returning \code{(false, nil)} for an untrusted,
incomplete or expired chain and \code{(false, err)} only for input it cannot
evaluate at all; and it never returns true by default --- with an empty root
pool nothing verifies.

The trusted root pool (\code{precompile/ai/tee\_roots.go}) currently holds
exactly one certificate: the Intel SGX Provisioning Certification Root CA
(SHA-256 fingerprint \code{44:A0:19:6B:\ldots:B6:74:D3}), which anchors Intel
SGX and TDX DCAP chains. The verification core is vendor-agnostic; a vendor
whose root is not embedded fails closed.

The \code{attestation} precompile at
\code{0x0300000000000000000000000000000000000001} is a stateless adapter over
that verifier, with selectors \code{0x01} \code{VerifyNVTrust}, \code{0x02}
\code{VerifyTPM}, \code{0x03} \code{VerifyCompute}, \code{0x04}
\code{CreateAttestation} and \code{0x05} \code{GetDeviceStatus}. It keeps no
device registry: a prior process-global map was a consensus-split hazard and was
removed, so \code{GetDeviceStatus} always reports not-found and a quote must be
presented to be believed. Attestation expiry uses the block timestamp plus a
3600-second TTL. Its \code{SupportedGPUModels} list is H100, H200, B100, B200,
GB200 and RTX PRO 6000.

\subsection{The compute-market precompile, and its admission defect}
\label{sec:defect}

The AI precompile range is \code{0x0300\ldots00} through \code{0x0300\ldots ff}
(\code{precompile/modules/registerer.go}). Within it:
\code{0x\ldots0000} AI mining, \code{0x\ldots0001} attestation,
\code{0x\ldots0002} model registry, \code{0x\ldots0003} inference,
\code{0x\ldots0004} AIVM bridge, and \code{0x\ldots0010} the compute market
(\code{precompile/compute/contract.go}, interface
\code{IComputeMarket.sol}).

The compute market implements five operations --- \code{RegisterProvider},
\code{SubmitJob}, \code{ClaimReward}, \code{VerifyCompute}, \code{GetPrice} ---
over a keccak-prefixed storage layout, with a job status of open, claimed or
verified. Its per-operation gas is 25{,}000 / 10{,}000 / 5{,}000 / 50{,}000 /
5{,}000, and \code{GetPrice} quotes $1000$ wei per input byte and refuses a
product wider than 256 bits rather than wrapping.

Its verification step, however, is not verification. \code{verifyCompute}
computes

\begin{center}
\code{expected = keccak256(jobID || outputHash)}
\end{center}

and marks the job verified when the caller's \code{attestation} argument equals
that value (\code{contract.go}, lines 270--283). Both operands are already
public: \code{jobID} is the argument, and \code{outputHash} was written by the
claimant at \code{ClaimReward}. Anyone can compute the expected value, so the
check establishes that the caller can hash --- not that a computation occurred,
not where, and not on what. It is a well-formedness checksum wearing the name of
an attestation check.

This is the single most consequential thing to fix, and the fix is not to write
a new verifier. §\ref{sec:changes} routes admission through the one verifier
that already exists.

\subsection{Model commitments}

The model registry precompile (\code{0x0300\ldots0002},
\code{precompile/modelregistry/registry.go}) holds a governance-adopted,
versioned commitment per model name:
$\code{GetApproved}(\textit{name}) \to (\textit{version}, \textit{weight})$,
where \textit{weight} is a commitment to the weights and a zero value means no
version has been adopted. Adoption is gated on an admin set held in the same
precompile's storage. This is the on-chain answer to ``which weights are
canonical'', and §\ref{sec:slash} uses it as the reference for substitution.

\subsection{What an operator can say about itself}
\label{sec:evidence}

\code{hanzod} produces a measurement and signed evidence about a microVM boot
(crate \code{hanzo-attest}, \code{hanzoai/node/hanzo-libs/hanzo-attest},
specified in \emph{Attested Compute in \code{hanzod}}). Its boot inputs are the
monitor binary digest, the operating-system asset version, the kernel digest,
the root-filesystem digest, the machine shape, and the guest argv; the
measurement is $\mathrm{BLAKE3}$ over a length-prefixed canonical encoding under
the domain tag \code{hanzo.boot.v1}. Evidence is $(M, \kappa, p, \sigma, t)$
signed over $\code{hanzo.evidence.v1} \| \kappa \| M \| t$, with kinds
$\kappa \in \{\code{host}=1, \code{snp}=2, \code{tdx}=3\}$.

Only \code{host} is implemented: the node signs its own measurement with its own
ed25519 identity key. \code{snp} and \code{tdx} are values with no attester and
no verifier. The distinction is the entire security argument. Host-asserted
evidence establishes that a named node made a claim and that the claim was not
altered; it does not establish that the node fed those files to the monitor,
because nothing outside the node observed the boot. A dishonest operator
measures the files it wishes to claim, signs, and boots whatever it likes --- and
the signature verifies.

The monitor is \code{hanzo-vm} (\code{hanzoai/vm}), which boots a k3s microVM
with published medians of 0.30\,s (Apple M4 Max, Virtualization.framework),
0.27\,s (Ryzen AI Max+ 395, cloud-hypervisor) and 0.30\,s (GB10, KVM). The
control-plane image reaching such a machine is \code{ghcr.io/hanzoai/cloud}.

\subsection{What the reference fleet can actually attest}
\label{sec:fleet}

\emph{Measured.} We inspected the two reference hosts.

\begin{table}[h]
\centering
\begin{tabular}{@{}lll@{}}
\toprule
 & \code{evo.local} & \code{spark.local} \\
\midrule
Architecture        & x86\_64                        & aarch64 \\
CPU                 & AMD Ryzen AI MAX+ 395          & --- \\
Cores               & 32                             & 20 \\
KVM                 & present                        & present \\
GPU                 & none reported by \code{nvidia-smi} & NVIDIA GB10 \\
\code{/dev/sev}     & absent                         & n/a \\
\code{kvm\_amd sev} & \code{N}                       & n/a \\
\code{\~{}/.hanzo/vm} assets & present               & --- \\
\bottomrule
\end{tabular}
\caption{Reference hosts, inspected directly. The two rows that matter for
admission are the SEV rows: SEV is not enabled on the AMD host, and the ARM
host is not an SEV or TDX platform at all.}
\end{table}

Two consequences follow, and neither is hypothetical:

\begin{enumerate}
  \item Neither host can produce a quote the implemented verifier would accept.
    The only embedded root is Intel SGX (§\ref{sec:verifier}); the AMD host has
    SEV disabled and exposes no \code{/dev/sev}; the ARM host is neither.
  \item GB10 is not in \code{SupportedGPUModels}, so even a GPU attestation
    path would report it as not confidential-computing capable.
\end{enumerate}

The design consequence is structural, not a stopgap: \textbf{the chain must be
able to price two trust bases, because a real fleet contains both}. A settlement
rule that admits only hardware-rooted evidence settles nothing on this fleet
today; a rule that treats host-asserted evidence as equivalent to a hardware
quote lies about what it checked. §\ref{sec:settle} resolves this by letting the
evidence kind modulate emission and nothing else.

\subsection{Where usage is metered today}

Per-request detail lands in the \code{hanzo.cloud\_usage} warehouse ledger
written by the AI plane (HIP-0106); the money lands on the org's commerce
ledger; \code{/v1/usage} (HIP-1313) composes reads over both and owns one series
of its own. HIP-1313's own motivation states the rule this proposal inherits:
answering ``what did this cost'' by copying three ledgers into a fourth store
produces a fourth number that can disagree with the ledger it came from. The
chain is not that fourth store.

\section{The chain}

\subsection{Placement}

\emph{Proposed.} Hanzo is a sovereign L1 whose parent is the Lux primary
network, with

\begin{align*}
  \code{security.Mode} \;=\; \bigl(\;
    &\code{RestakeParent} = \text{false},\;
     \code{Admission} = \code{Open},\\
    &\code{Threshold} = T,\;
     \code{Manager} = \code{Contract}\;\bigr)
\end{align*}

and \code{ManagerChainID} naming the Hanzo EVM chain, \code{ManagerAddress} the
validator-manager contract on it. \code{Mode.Sovereign()} holds. The genesis
validator list is carried by the same \code{CreateNetworkTx}
(§\ref{sec:sovereign}), and the chain itself is created by a \code{CreateChainTx}
in the same block.

Following the workspace convention that a sovereign L1's primary network
identifier equals its EVM chain identifier, Hanzo's identifiers are those
already recorded in HIP-0024: 36963 on mainnet and 36962 on testnet, ticker AI.
One identifier per environment removes the wallet and tooling ambiguity of a
network identifier that differs from the chain identifier.

The chain inherits Lux's consensus, its post-quantum signing paths, its
precompile block, and its Warp messaging. It does not inherit the primary
network's validator set --- that is what sovereign means --- and it does not
depend on the C-Chain, which is the Lux primary network's own EVM and is used by
no other network.

HIP-0024 records the engine as \code{nova} with sampling parameters
$k = 20$, $\alpha = 15$. Nova is the linear-chain driver
(\code{consensus/protocol/nova}) and Quasar is the post-quantum finality layer
that composes with such drivers (§\ref{sec:consensus}); the two names describe
different layers and are not alternatives. Nothing in this proposal depends on
which driver or which of Quasar's signing modes is configured, for the reason
given in §\ref{sec:consensus}: the settlement rules need deterministic finality and a validator
weight the P-Chain can be told to change, and both hold under every mode.

\subsection{Two sets, deliberately distinct}

The validator set secures the ledger. The operator set serves compute. They are
different populations with different capital, different hardware and different
failure modes, and LP-323 is an argument for keeping capacity, work and identity
apart. This proposal keeps them apart:

\begin{itemize}
  \item A \emph{validator} holds a P-Chain L1 validation with a
    \code{ValidationID}, a BLS key, and a continuous-fee balance topped up by
    \code{IncreaseL1ValidatorBalanceTx}.
  \item An \emph{operator} holds an EVM-level bond in the settlement contract
    and an ed25519 identity that signs evidence.
\end{itemize}

A principal may be both. Nothing requires it, and §\ref{sec:slash} is careful
about which of the two a slash touches.

\section{The ledger}
\label{sec:ledger}

\subsection{One asset, three accounts}

\emph{Proposed.} There is one asset, AI. There is no separate credit token. A
non-transferable credit unit would create a second price, and a second price
needs a reconciliation, and the reconciliation is the defect HIP-1313 exists to
avoid. Credit is not a different asset; it is AI under a spend authority.

For each buyer $b$ and operator $o$ the chain holds:

\begin{align*}
  \esc(b)  &\in \mathbb{N} && \text{AI deposited by $b$, spendable only through leases} \\
  \res(b)  &\in \mathbb{N} && \text{the part of $\esc(b)$ committed to open leases} \\
  \bond(o) &\in \mathbb{N} && \text{AI locked by $o$, slashable}
\end{align*}

with the invariant $\res(b) \le \esc(b)$ maintained at every transition.

\begin{definition}[Credit operations]
\code{deposit}$(v)$ adds $v$ to $\esc(b)$.
\code{withdraw}$(v)$ requires $v \le \esc(b) - \res(b)$ and subtracts $v$.
Nothing else moves $\esc$ except settlement.
\end{definition}

Three operations are the whole credit ledger. Money enters once per deposit, not
once per request; §\ref{sec:split} explains why the per-request meter must stay
off the chain.

\subsection{Leases}

\begin{definition}[Lease]
A \emph{lease} is
$\ell = (b,\, o,\, m,\, \nu,\, \pi,\, \bar u,\, [e_1, e_2])$: a buyer, an
operator, a model name $m$ in the registry, a \emph{unit} $\nu$, a price $\pi$ in
AI per unit, a cap $\bar u$ on units, and a window of epochs. Its
\emph{obligation} is $\pi \bar u$.
\end{definition}

$\nu$ is an opaque identifier and the chain never interprets it. Whether a unit
is an output token, a GPU-second, a request or a frame is a question the two
counterparties answer between themselves and the meter records off the chain
(§\ref{sec:split}); the chain settles a scalar they both signed for. Putting a
unit taxonomy on the chain would fix in consensus a vocabulary that changes with
every new modality, and would give the chain a second opinion about work it
cannot observe.

Opening a lease reserves the buyer's side and charges the operator's book:

\begin{norm}[Opening]
\label{rule:open}
$\code{open}(\ell)$ MUST fail unless
$\res(b) + \pi \bar u \le \esc(b)$ and
\[
  \lambda \Bigl(\pi_\ell \bar u_\ell +
    \sum_{\ell' \in \open(o)} \pi_{\ell'} \bar u_{\ell'}\Bigr) \;\le\; \bond(o),
\]
where $\open(o)$ is $o$'s set of live leases and $\lambda \in (0,1]$ is the
slash fraction of §\ref{sec:slash}. On success $\res(b) \mathrel{+}= \pi\bar u$
and $\ell$ joins $\open(o)$.
\end{norm}

The second condition is the invariant that makes slashing sound. It says the
bond covers the book: an operator may take on obligations only up to the point
where the maximum slash they can incur is still collectible. Most designs assert
a slash percentage and never check that the stake can pay it; this checks it at
the only moment when checking is cheap.

\section{Claims and admission}
\label{sec:admit}

\subsection{The one object that crosses}

\emph{Proposed.} Work is served continuously and settled periodically. Let an
\emph{epoch} $e$ be a fixed number of blocks. Within an epoch an operator
accumulates \emph{receipts} off the chain; at the epoch's close it submits one
\emph{claim} on the chain.

\begin{definition}[Receipt]
A receipt is $r = (\ell,\, u_r,\, h_r,\, t_r)$: the lease it serves, the units of
work, a digest of the request and response bodies, and a time. Receipts never
appear on the chain.
\end{definition}

\begin{definition}[Claim]
A claim is
\[
  C \;=\; \bigl(o,\; e,\; \rho,\; M,\; \kappa,\; Q,\;
    \{(\ell_i, u_i)\}_{i=1}^{n},\; \sigma_o \bigr)
\]
where $\rho$ is the Merkle root over the epoch's receipts, $M$ is the boot
measurement of §\ref{sec:evidence}, $\kappa$ the evidence kind, $Q$ the
attestation receipt and signature, $(\ell_i, u_i)$ the per-lease unit totals, and
$\sigma_o$ the operator's signature over the whole.
\end{definition}

A claim is $O(1)$ in the number of requests it covers. Its size is
$32 + 32 + 1 + |Q| + 40n + |\sigma_o|$ bytes, dominated by $Q$, whose bulk is a
DER certificate chain. At one claim per operator per epoch, chain load is
$O(|\mathcal{O}|)$ and independent of traffic: $10^4$ operators on hourly epochs
is under three claims per second. This is the reason for the epoch structure, and
it is the whole content of the on-chain/off-chain boundary.

\subsection{Admission}

\begin{definition}[Admission]
\label{def:adm}
$\Adm(C)$ holds exactly when all of the following hold. Each is a total,
deterministic function of chain state and $C$.

\begin{enumerate}[label=A\arabic*.,leftmargin=2.4em]
  \item \textbf{Evidence verifies.}
    $\code{ai.VerifyTEE}(Q.\text{receipt}, Q.\sigma) = \text{true}$ for
    $\kappa \in \{\code{snp}, \code{tdx}\}$, against the embedded root pool at
    the report's own timestamp (§\ref{sec:verifier}); or, for
    $\kappa = \code{host}$, $\sigma$ verifies under the operator's registered
    ed25519 identity over
    $\code{hanzo.evidence.v1} \,\|\, \kappa \,\|\, M \,\|\, t$.
  \item \textbf{Evidence binds this claim.} The report's 8-byte nonce equals the
    low 8 bytes of
    $H(\code{hanzo.claim.v1} \,\|\, o \,\|\, e \,\|\, \rho \,\|\, M)$.
  \item \textbf{Evidence is fresh.} The report's timestamp lies within epoch
    $e$'s window.
  \item \textbf{Device is the operator's.} The report's 32-byte deviceID is one
    registered to $o$.
  \item \textbf{Unspent.} The work identifier $\omega = H(o \,\|\, e)$ is not in
    the spent set. On admission it is marked spent.
  \item \textbf{Bonded.} $\bond(o) \ge B_{\min}$ and $o$ is not in an unbonding
    period.
  \item \textbf{Arithmetic.} Every $\ell_i$ names $o$, is live in $e$, and
    satisfies $u_i + (\text{units already settled on } \ell_i) \le \bar u_{\ell_i}$.
  \item \textbf{Model is adopted.} For each distinct $m$ among the $\ell_i$,
    $\code{GetApproved}(m)$ returns a nonzero weight commitment.
\end{enumerate}
\end{definition}

A2 is the anti-replay binding, and it uses a field that already exists. The
48-byte report header carries an 8-byte nonce whose purpose in the verifier is
exactly this and which the verifier does not otherwise constrain. Binding it to
$(o, e, \rho, M)$ means a quote produced for one operator's epoch cannot be
presented for another's: changing any of the four changes the nonce, and the
nonce is inside the signed header. Without A2, an operator with no TEE replays a
competitor's valid quote.

$M$ belongs inside that hash for a reason worth stating, because the alternative
looks equally correct and is not. The Lux receipt header is
$(\text{deviceID}, \text{timestamp}, \text{nonce})$ and has no measurement field,
so the vendor-rooted signature covers the measurement only if the measurement is
folded into one of the three. The nonce is the only one that is free. With $M$
outside the hash, $M$ would be covered by $\sigma_o$ alone --- the operator's own
signature --- and a hardware quote would attest to a device and a time while the
statement about what booted stayed self-asserted. Folding $M$ in makes the
vendor-rooted signature cover it transitively, which is what
$\mu(\code{snp}) > \mu(\code{host})$ is paying for.

The cost of that choice is the nonce's width. Eight bytes gives a forger 64 bits
of second-preimage work to find a different $(o, e, \rho, M)$ matching a nonce
they already hold a valid quote for, and they must additionally control an
operator the deviceID is registered to (A4). Sixty-four bits is adequate against
a targeted forgery and is not a birthday setting, since the target nonce is
fixed. It is nonetheless the narrowest quantity in this design, and widening the
receipt's nonce field is the change that would widen it.

A5 is the second replay guard and covers the case where an operator submits the
same claim twice. The pattern is implemented in the AI mining precompile:
\code{IsSpent}/\code{MarkSpent} over a
$\mathrm{BLAKE3}(\text{prefix} \| \omega)$ key
(\code{precompile/ai/ai\_mining.go}).

\subsection{What admission does not establish}

Stating this precisely is part of the specification, not a caveat appended to it.

\begin{itemize}
  \item It does not establish that the adopted model produced the outputs. It
    establishes that the operator's declared boot inputs hash to $M$, and, under
    $\kappa \in \{\code{snp},\code{tdx}\}$, that a genuine platform signed $M$.
    Binding $M$ to the registry's weight commitment is §\ref{sec:slash}'s
    substitution rule and is a claim about \emph{declared} inputs.
  \item It does not establish result quality. No hash of an output is evidence
    that the output is right.
  \item Under $\kappa = \code{host}$ it does not establish that the measured
    files were the files the monitor was given, because nothing outside the node
    observed the boot. Host-asserted evidence is worth the reputation of the
    signer, and §\ref{sec:settle} prices it accordingly rather than pretending
    otherwise.
  \item It does not establish that the units $u_i$ are the units the buyer
    actually consumed. It establishes that the operator committed to a receipt
    set whose root is $\rho$ and signed for the totals. Disagreement is a
    dispute, and disputes are HIP-0095's subject; the buyer holds its own
    receipts and opens a Merkle path against $\rho$.
\end{itemize}

\section{Settlement}
\label{sec:settle}

\emph{Proposed.} At the close of epoch $e$, for every admitted claim, two
independent terms are computed. Keeping them independent is the point: one is
conserved and one is minted, and a design that sums them into a single ``reward''
loses the ability to take emission to zero without breaking pay-per-use.

\subsection{The conserved term}

For admitted claim $C$ and each $(\ell_i, u_i)$, with $b_i$ the lease's buyer:

\begin{align*}
  R_i     &= \pi_{\ell_i}\, u_i && \text{revenue on this lease} \\
  \esc(b_i) &\mathrel{-}= R_i,\qquad \res(b_i) \mathrel{-}= R_i \\
  \pay(o) &\mathrel{+}= (1-\phi)\,R_i && \text{operator's share} \\
  \text{treasury} &\mathrel{+}= \phi\, R_i && \text{protocol fee}
\end{align*}

where $\phi \in (0,1)$ is the protocol fee rate. Nothing is minted here. When a
lease closes with $u < \bar u$, the unspent reserve is released:
$\res(b) \mathrel{-}= \pi(\bar u - u)$.

Let $w_o = \sum_i R_i$ be the operator's admitted revenue in the epoch and
$W_e = \sum_{o} w_o$ the epoch's total.

\subsection{The minted term}

\begin{definition}[Emission]
\label{def:emit}
With epoch pool $E_e$, rate cap $\varepsilon > 0$, the evidence kind $\kappa$
carried by the operator's admitted claim, and multiplier
$\mu(\kappa) \in (0,1]$:
\[
  \emit(o) \;=\; \mu(\kappa)\cdot
    \min\!\left( E_e \cdot \frac{w_o}{W_e},\;\; \varepsilon\, w_o \right).
\]
Any part of $E_e$ not allocated by this formula is not minted.
\end{definition}

Two choices in that formula carry the design, and each replaces a mechanism that
would otherwise need its own section.

\paragraph{Weighting by paid work rather than by capacity.} A capacity
normalisation --- HIP-0096's Hanzo Compute Units, a weighted sum over GPU FLOPs,
VRAM, cores, memory bandwidth, storage IOPS and network --- is a number the
operator reports about itself. Weighting emission by it creates a direct
incentive to inflate it, which is why designs built on it need an anti-gaming
section chasing the inflation. Weighting by $w_o$ removes the incentive rather
than policing it: capacity nobody paid for earns nothing, and the number in the
weight is one a counterparty agreed to pay.

The cost of this choice is real and should be stated: emission follows price, so
an operator serving cheap models in a cheap region earns less emission per unit
of physical work than one serving expensive models. That is a consequence of
paying for value rather than for watts, and the alternative reintroduces
self-reported capacity.

\paragraph{Rate-capping rather than pooling.} A pure pro-rata split of a fixed
pool has a hole that is easy to miss: as $W_e \to 0$ the emission \emph{per unit
of revenue} diverges, so a single participant in a thin epoch collects nearly all
of $E_e$ for an arbitrarily small $w_o$. The $\varepsilon w_o$ term caps marginal
emission at $\varepsilon$ per unit of revenue regardless of how thin the epoch
is. It is one \code{min} and it closes the empty-epoch attack.

\paragraph{The multiplier is on emission only.} $\mu(\code{host}) <
\mu(\code{snp}) = \mu(\code{tdx})$. The protocol subsidises better evidence; it
does not reprice a contract the counterparties agreed. A buyer who paid $\pi$ per
unit gets charged $\pi$ per unit whatever the evidence kind, and the operator
receives $(1-\phi)R$ whatever the evidence kind. §\ref{sec:fleet} is why this
parameter exists at all: on the reference fleet today, $\kappa = \code{host}$ is
the only kind available, and a chain that could not settle it would settle
nothing.

\subsection{Collusion is unprofitable}

\begin{proposition}
\label{prop:wash}
Let a single economic principal control both a buyer $b$ and an operator $o$,
and route revenue $R$ through leases between them. Under
Definition~\ref{def:emit}, the principal's net change in AI holdings is at most
$(\varepsilon - \phi) R$. If $\varepsilon < \phi$ it is strictly negative for
every $R > 0$.
\end{proposition}

\begin{proof}
The conserved term moves $R$ from $\esc(b)$ to $\pay(o)$ less the fee: since $b$
and $o$ are the same principal, the internal transfer nets to zero and the
principal's holdings change by $-\phi R$. The minted term is bounded above by
$\mu(\kappa)\,\varepsilon\,R \le \varepsilon R$, because the \code{min} in
Definition~\ref{def:emit} is at most its second argument and $\mu \le 1$. The
sum is at most $(\varepsilon - \phi)R$, which is negative for $R > 0$ exactly
when $\varepsilon < \phi$.
\end{proof}

\begin{remark}
The bound holds for every epoch composition, including an epoch in which the
colluding pair is the only participant. That is what the rate cap buys; under a
pure pooled split the second term would be $E_e$ rather than $\varepsilon R$ and
the proposition would be false for small $R$.
\end{remark}

\begin{remark}
Proposition~\ref{prop:wash} bounds a principal who controls both sides. It says
nothing about an operator who serves real buyers badly; that is quality, and
quality is HIP-0095's subject. It also assumes the fee is not rebated to the
payer by some other mechanism --- a rebate is an $\varepsilon$ in disguise and
must be counted in $\varepsilon$.
\end{remark}

\begin{proposition}[Emission is bounded by paid work]
$\sum_o \emit(o) \le \min(E_e,\; \varepsilon W_e)$.
\end{proposition}

\begin{proof}
Immediate from Definition~\ref{def:emit}: term-wise, $\sum_o \mu \cdot \min(\cdot)
\le \sum_o E_e w_o / W_e = E_e$ and $\sum_o \mu \cdot \min(\cdot) \le
\sum_o \varepsilon w_o = \varepsilon W_e$.
\end{proof}

This is HIP-0096's stated invariant ``total rewards $\le$ total value created''
made structural rather than aspirational: with no paid work there is no
emission, and the unallocated pool is not minted.

\section{Slashing}
\label{sec:slash}

\subsection{Rejection is not slashing}

An operator whose attestation does not verify has produced nothing the chain can
admit. That is a rejection: the claim does not settle, and the operator has
wasted the gas. It is not, by itself, a slash. Slashing a failed verification
punishes an operator for a network fault, a clock skew, or an expired
certificate, and an operator who can be slashed for a transient failure will
withdraw its bond.

What is slashable is an \emph{obligation taken and not met}. The operator
accepted a lease under Rule~\ref{rule:open}, and the buyer's money was reserved
against it. If no admitted claim covers that lease by its deadline, the operator
took a commitment and did not deliver, and the reason --- no attestation, an
invalid attestation, a crashed node, a decision not to serve --- does not change
the fact.

This is one rule where a design that slashes both ``invalid'' and ``missing''
needs two, and it is strictly stronger than either alone: an invalid attestation
on a live lease produces no admitted claim, so non-delivery fires anyway.

\subsection{The three slashable events}

\emph{Proposed.} All three are decidable from chain state alone.

\begin{norm}[Non-delivery]
\label{rule:nd}
If lease $\ell \in \open(o)$ reaches epoch $e_2 + \Delta$ with settled units
$u < \underline{u}\,\bar u$ for a delivery floor $\underline{u} \in [0,1]$, then
$s = \min\bigl(\bond(o),\, \lambda\, \pi_\ell (\bar u_\ell - u)\bigr)$ is taken
from $\bond(o)$, the buyer's remaining reserve on $\ell$ is released, and $s$ is
credited to $\esc(b)$.
\end{norm}

Two details in that rule are load-bearing.

The slash is proportional to the obligation, not a flat fraction of the bond. A
flat fraction under-collateralises a large book and over-punishes a small one;
proportionality makes the required bond a function of the book the operator chose
to take, which is what Rule~\ref{rule:open} enforces at the moment of taking it.

The slash goes to the buyer, not to the treasury, and it is measured on the
\emph{undelivered} part $\bar u_\ell - u$. The buyer is the party that was
harmed: it reserved capital against a promise and got a fraction of it. Sending
the slash to the treasury would make the protocol the beneficiary of its own
operators' failures, which is an incentive nobody should want to build. Crediting
$\esc(b)$ rather than paying out keeps the money inside the credit ledger, where
the buyer can spend it with a different operator without a second transaction.

Together, Rules \ref{rule:open} and \ref{rule:nd} give:

\begin{proposition}[Slashes are collectible]
Under Rule~\ref{rule:open}, at any time
$\sum_{\ell \in \open(o)} \lambda \pi_\ell \bar u_\ell \le \bond(o)$. Each
non-delivery slash is $\lambda \pi_\ell (\bar u_\ell - u) \le
\lambda \pi_\ell \bar u_\ell$, so the sum of all non-delivery slashes an operator
can simultaneously incur does not exceed its bond, and the $\min$ in
Rule~\ref{rule:nd} never binds.
\end{proposition}

\begin{norm}[Equivocation]
If two admitted claims $C, C'$ of the same operator cover overlapping intervals
of the same lease, $\lambda_{\mathrm{eq}} \bond(o)$ is taken from $\bond(o)$ and
both claims' settlements are reversed.
\end{norm}

\begin{norm}[Substitution]
If an admitted claim names model $m$ and the measurement $M$ in its evidence is
not the measurement the operator registered against $(m, \text{version})$ for the
weight commitment $\code{GetApproved}(m)$ returns, then
$\lambda_{\mathrm{sub}} \bond(o)$ is taken from $\bond(o)$ and the claim's
settlement is reversed.
\end{norm}

These two are fractions of the bond rather than of an obligation, because
neither harm is bounded by one lease: an equivocating or substituting operator
has misreported to every buyer it served that epoch. Their proceeds are split
pro-rata to $\esc(b)$ across the affected buyers and the remainder goes to the
treasury, which is the only case in this design where the treasury receives a
slash --- and it receives only what no identifiable buyer was owed.

\begin{remark}
Substitution is a statement about \emph{declared} inputs. Under
$\kappa = \code{host}$ an operator who intends to substitute simply declares the
measurement it is supposed to have and boots something else; the rule catches
mistakes and lazy operators, not determined ones. Under
$\kappa \in \{\code{snp},\code{tdx}\}$ the measurement is signed by a platform
the operator does not control and the rule binds. This is precisely the gap
$\mu(\kappa)$ prices, and it is the reason a hardware attester is the highest-value
piece of work outstanding in the fleet rather than a refinement.
\end{remark}

\subsection{Reaching the validator set}

A slash lands on the EVM bond, which is where operators keep capital. If the
slashed operator is also an L1 validator, the validator-manager contract named by
\code{ManagerAddress} (§\ref{sec:sovereign}) reduces its weight by emitting an
\code{L1ValidatorWeight} Warp payload to the Lux P-Chain carrying the
\code{ValidationID} and the new weight; removal is the reserved nonce
$2^{64}-1$ with weight $0$. That path is implemented and specified in LP-77; this
proposal adds no new cross-chain mechanism.

An operator that is not a validator has no validator weight to reduce, and
nothing about its slash touches consensus. Keeping the two sets distinct
(§\ref{sec:sovereign}) is what makes that sentence true.

\section{On chain and off chain}
\label{sec:split}

The boundary is one sentence: \textbf{the chain holds commitments and verifies
signatures; the fleet holds content and produces signatures; exactly one object
crosses, and it is the claim.}

\begin{longtable}{@{}p{0.30\textwidth}p{0.16\textwidth}p{0.46\textwidth}@{}}
\toprule
Object & Lives & Why \\
\midrule
\endhead
Credit balance, reserve & chain & It is the money. \\
Operator bond, its lock & chain & It is the collateral a slash collects from. \\
Lease terms $(b,o,m,\nu,\pi,\bar u,\text{window})$ & chain & Both sides must be
held to the same numbers, and Rule~\ref{rule:open} is checked against them. The
unit $\nu$ is stored and never interpreted. \\
Admitted-claim record & chain & Equivocation and non-delivery are decided by
comparing claims. \\
Spent set over $\omega$ & chain & Replay is only detectable against a durable set. \\
Model commitment & chain & Substitution needs a reference the operator does not
control. \\
Emission parameters, epoch pool & chain & They determine minting. \\
\midrule
Prompts, inputs, outputs & fleet & Never on chain, in any form. Confidentiality
is not a property a public ledger can hold, and LP-324 is explicit that it is
not a receipt property either. \\
Per-request meter & fleet & \code{hanzo.cloud\_usage} (HIP-0106),
\code{/v1/usage} (HIP-1313). Copying it on chain creates a second number to
reconcile. \\
Receipts and the tree over them & fleet & $O(\text{requests})$; the root is
$O(1)$ and is what the chain needs. \\
Routing, scheduling, placement & fleet & No settlement rule reads them. \\
Quote generation & fleet & The platform signs; the chain verifies. \\
Price discovery & fleet & The chain records the price the lease agreed, not how
it was found. \\
\midrule
\textbf{The claim} & \textbf{crosses} & $(o,e,\rho,M,\kappa,Q,\{(\ell_i,u_i)\},\sigma_o)$:
one per operator per epoch, $O(1)$ in requests. \\
\bottomrule
\caption{The split. A row is on the chain when a settlement rule reads it and no
counterparty can be trusted to hold it; otherwise it is on the fleet.}
\end{longtable}

The test each row answers is not ``is it important'' but ``does a settlement rule
in this document read it''. Nothing else earns chain state. Applying that test is
what keeps the claim $O(1)$: the chain reads $\rho$, never a receipt, so traffic
does not enter the state.

\subsection{The dispute path}

A buyer disagreeing with $u_i$ holds its own receipts and the operator's $\rho$.
It opens a challenge (HIP-0095) carrying a Merkle path. The chain adjudicates a
path against a root, which is a hash comparison, not a re-execution. This
proposal specifies no more of the challenge protocol than the commitment it needs
--- $\rho$, in the claim --- because HIP-0095 already specifies the rest, and two
specifications of one protocol is the defect to avoid.

\section{Transaction and precompile surface}

\emph{Proposed.} The settlement logic is an EVM contract on the Hanzo chain
rather than a new precompile. Precompiles are for primitives that cannot be
expressed in the EVM at reasonable cost --- signature verification, lattice
arithmetic, hashing --- and the AI range already holds the ones this design uses.
Settlement is bookkeeping over those primitives, it changes with governance, and
a contract can be upgraded on the chain's own schedule while a precompile change
is a network upgrade.

\begin{longtable}{@{}p{0.30\textwidth}p{0.22\textwidth}p{0.40\textwidth}@{}}
\toprule
Call & Caller & Effect \\
\midrule
\endhead
\code{deposit()} payable & buyer & $\esc(b) \mathrel{+}= \text{value}$ \\
\code{withdraw(v)} & buyer & requires $v \le \esc(b)-\res(b)$ \\
\code{bond()} payable & operator & $\bond(o) \mathrel{+}= \text{value}$ \\
\code{unbond(v)} & operator & requires
$\bond(o) - v \ge \lambda \sum_{\ell \in \open(o)} \pi_\ell \bar u_\ell$;
withdrawable at $\max_{\ell \in \open(o)} (e_2(\ell) + \Delta)$ \\
\code{device(id, quote)} & operator & registers a deviceID for A4 \\
\code{open(o,m,$\nu$,$\pi$,$\bar u$,window)} & buyer or its delegate & Rule~\ref{rule:open} \\
\code{claim(C)} & operator & Definition~\ref{def:adm}; on success records and
marks $\omega$ spent \\
\code{settle(e)} & anyone & §\ref{sec:settle} over epoch $e$'s admitted claims \\
\code{report(event, subject)} & anyone & §\ref{sec:slash}; \emph{subject} is a
lease or a pair of claims \\
\bottomrule
\caption{The surface. \code{settle} is permissionless because it is a pure
function of admitted claims; whoever pays its gas cannot change its result.}
\end{longtable}

The precompiles it calls, all existing: \code{0x0300\ldots0001} for A1 under
$\kappa \in \{\code{snp},\code{tdx}\}$, and \code{0x0300\ldots0002} for A8 and
the substitution rule.

\code{unbond} does not wait for the operator's book to empty --- an operator
with a rolling book would then never be able to leave --- but it does hold the
funds until every live obligation has passed its non-delivery deadline. The two
conditions are the exit-side counterpart of Rule~\ref{rule:open}: the first keeps
the remaining bond sufficient for the remaining book, and the second keeps it
present long enough for Rule~\ref{rule:nd} to collect.

\code{device} takes a quote rather than a bare identifier, so a deviceID is
registered by a party that can make that device sign. Registering someone else's
deviceID gains nothing regardless, because A2 requires a nonce that binds the
claiming operator and only the device can produce one; the possession check keeps
the registry honest rather than doing security work of its own.

\code{open} accepts a delegate so that the gateway can open leases under a
spend authority the buyer granted, which is how a pay-per-use API becomes a lease
without the buyer signing per request. The authority's own shape --- its bound,
its expiry, its revocation --- is HIP-0025 and LP-325, and is not restated here.

\section{Parameters}

\emph{Proposed.} Values are initial proposals subject to governance; the
constraints between them are not.

\begin{longtable}{@{}llp{0.46\textwidth}@{}}
\toprule
Symbol & Initial & Meaning and constraint \\
\midrule
\endhead
$L$ & 1 hour of blocks & Epoch length. Longer amortises $|Q|$ over more work;
shorter shortens the buyer's exposure to an operator that stops. \\
$\phi$ & $0.05$ & Protocol fee on conserved revenue. \\
$\varepsilon$ & $0.02$ & Emission per unit revenue. MUST satisfy
$\varepsilon < \phi$ (Proposition~\ref{prop:wash}). \\
$E_e$ & schedule & Epoch pool. An upper bound only; the unallocated part is not
minted. \\
$\mu(\code{host})$ & $0.25$ & Emission multiplier for host-asserted evidence. \\
$\mu(\code{snp}), \mu(\code{tdx})$ & $1.0$ & Multiplier for hardware-rooted
evidence. \\
$\lambda$ & $1.0$ & Non-delivery slash as a fraction of the \emph{undelivered}
obligation, paid to the buyer. $\lambda = 1$ makes a defaulting operator forfeit
the full value of what it promised and did not serve. Also the coefficient in
Rule~\ref{rule:open}, so raising it tightens every operator's book. \\
$\lambda_{\mathrm{eq}}$ & $0.10$ & Equivocation slash, fraction of bond. \\
$\lambda_{\mathrm{sub}}$ & $0.10$ & Substitution slash, fraction of bond. \\
$\underline{u}$ & $0.90$ & Delivery floor before non-delivery fires. \\
$B_{\min}$ & governance & Minimum bond to submit any claim. Distinct from the
validator minimum stake, which HIP-0024 records as 2000 AI. \\
$\Delta$ & 1 epoch & Grace after a lease's window before non-delivery fires. \\
\bottomrule
\end{longtable}

\section{Adversaries}

\begin{description}[leftmargin=1.4em]
  \item[A network adversary] between an operator and the chain can modify what is
    in flight. A claim is self-contained and signed, so this fails: the chain
    recomputes the nonce binding from $(o,e,\rho,M)$ and checks the signature
    without trusting the transport.
  \item[An operator without a TEE] wants a hardware-rooted multiplier. Replaying
    another operator's quote fails A2, because the nonce commits to $o$ and is
    inside the signed report header. Fabricating a quote fails A1, because the
    chain does not terminate at an embedded root.
  \item[An operator resubmitting] a settled epoch fails A5.
  \item[A dishonest operator with host-asserted evidence] declares the
    measurement it is supposed to have and runs something cheaper. Nothing in
    this document detects that, and §\ref{sec:evidence} explains why nothing
    can: the root of trust is the party whose honesty is in question.
    $\mu(\code{host}) < 1$ prices it; only a hardware attester removes it.
  \item[A buyer and operator under one principal] farming emission is bounded by
    Proposition~\ref{prop:wash} and loses money whenever $\varepsilon < \phi$.
  \item[An operator overstating units] $u_i$ within its caps is not caught by
    admission --- admission checks arithmetic against caps, not against reality.
    It is caught by the buyer, who holds receipts and opens a Merkle path
    (HIP-0095). The chain's contribution is the commitment $\rho$, without which
    the dispute has nothing to be about.
  \item[An operator taking more obligations than it can serve] is refused at
    Rule~\ref{rule:open}, which is the only moment the check is cheap.
  \item[A validator equivocating on consensus] is out of scope here and is
    Quasar's subject; this document's slashes are about compute obligations, not
    about block production.
\end{description}

\section{Changes required in existing code}
\label{sec:changes}

\emph{Proposed}, and deliberately small. Every item removes something or routes
an existing call; none adds a new cryptographic primitive.

\begin{enumerate}[leftmargin=1.8em]
  \item \textbf{The compute precompile's \code{verifyCompute} must not be the
    admission check.} As implemented (§\ref{sec:defect}) it compares the caller's
    argument to \code{keccak256(jobID || outputHash)}, which the caller can
    compute. Admission per Definition~\ref{def:adm} routes A1 through
    \code{0x0300\ldots0001} (which is a stateless adapter over
    \code{ai.VerifyTEE}) and never through this function. The precompile's
    \code{verifyCompute} should either take a real quote and delegate, or be
    removed; keeping a function named \code{verifyCompute} that verifies a hash
    of public inputs is the kind of name that gets trusted by the next reader.

  \item \textbf{Embed the vendor roots the fleet actually uses.}
    \code{tee\_roots.go} holds one root, Intel SGX. §\ref{sec:fleet} measured a
    fleet on which no host can chain to it. AMD SEV-SNP and NVIDIA device
    identity roots must be appended for $\kappa \in \{\code{snp}\}$ and any GPU
    path to be reachable; the verification core is already vendor-agnostic and
    needs no change.

  \item \textbf{Add GB10 to \code{SupportedGPUModels}, or state why not.} The
    list carries GB200 and not GB10 while the reference ARM host is a GB10. One
    of the two is wrong and the discrepancy should be resolved deliberately.

  \item \textbf{Fix the AI mining spent-set address.} \code{module.go} registers
    the AI mining precompile at
    \code{0x0300000000000000000000000000000000000000}, while
    \code{ai\_mining.go}'s \code{precompileAddr} is
    \code{0x0000000000000000000000000000000000000300} --- the same nibbles, the
    other end of the word. The spent set is therefore written under an account
    the module does not own. A5 depends on this set, so it must be one address.

  \item \textbf{Attach an attester for \code{Kind::Snp} and \code{Kind::Tdx} in
    \code{hanzo-attest}.} The evidence structure needs no change; the kind becomes
    \code{snp} or \code{tdx} and the signer becomes the vendor chain, which item 2
    makes verifiable. The platform's user-data field (\code{REPORT\_DATA} on
    SEV-SNP, \code{REPORTDATA} on TDX) carries different content at two different
    times, and conflating them is easy: a quote taken \emph{at boot} has no claim
    to bind and carries $M$, which is what the \code{hanzod} paper proposes; a
    quote taken \emph{at epoch close} carries
    $H(\code{hanzo.claim.v1}\|o\|e\|\rho\|M)$, which is A2 and which already
    covers $M$. Settlement needs the second. This is the item that moves an
    operator from $\mu(\code{host})$ to $\mu(\code{snp})$, and it is the
    highest-value one on this list.

  \item \textbf{Give \code{hanzo-attest} a durable registry.} The implemented
    registry is the node's own in-process view and does not outlive the process.
    A claim is computed over an epoch's history, so the record must survive the
    epoch. The record is already self-contained and signed, so the transport need
    not be trusted; what durability adds is history.
\end{enumerate}

\section{Limitations}

\begin{enumerate}[leftmargin=1.8em]
  \item Under $\kappa = \code{host}$ --- the only kind implemented, and the only
    kind the reference fleet can produce --- admission establishes intent, not
    execution. Every claim settled this way is worth the operator's reputation.
    The design's response is to price it, not to hide it.
  \item Admission checks unit arithmetic against caps, not against delivered
    work. Overstatement within caps is a dispute, not a rejection, and depends on
    buyers actually disputing.
  \item The substitution rule binds a declared measurement to a registry
    commitment. It does not establish that the committed weights were loaded.
    LP-313 and LP-324's model-identity witness are the constructions that would.
  \item Emission weighted by revenue pays less for cheap work than a
    capacity-weighted scheme would, by construction.
  \item Proposition~\ref{prop:wash} assumes the fee is not rebated. Any mechanism
    that returns fee revenue to payers must be counted inside $\varepsilon$ or
    the bound is void.
  \item The epoch structure delays an operator's revenue by up to one epoch and
    delays detection of an operator that has stopped by up to $L + \Delta$.
    Shortening $L$ costs more on-chain verifications of $Q$.
  \item A claim's size is dominated by the certificate chain in $Q$. A fleet of
    $10^4$ operators on hourly epochs is under three claims per second, but the
    verification cost per claim is a chain-path validation, not a hash.
\end{enumerate}

\section{References}

\begin{description}[leftmargin=0em,labelsep=0.6em]
  \item[Lux, implementation.]
    \code{node/vms/platformvm/security/mode.go};
    \code{node/vms/platformvm/txs/\{create\_network\_tx, convert\_network\_tx,
    create\_chain\_tx, register\_l1\_validator\_tx,
    increase\_l1\_validator\_balance\_tx, disable\_l1\_validator\_tx\}.go};
    \code{node/vms/platformvm/warp/message/l1\_validator\_weight.go};
    \code{node/vms/platformvm/network/zap\_native\_admission.go};
    \code{consensus/protocol/quasar};
    \code{precompile/ai/\{tee\_attest, tee\_roots, ai\_mining\}.go};
    \code{precompile/attestation};
    \code{precompile/modelregistry};
    \code{precompile/compute};
    \code{precompile/modules/registerer.go}.
  \item[Hanzo, implementation.]
    \code{hanzoai/node/hanzo-libs/hanzo-attest};
    \code{hanzoai/vm};
    \code{hanzoai/cloud} \code{apps/usage}.
  \item[Hanzo proposals.]
    HIP-0024 sovereign L1 chain architecture;
    HIP-0025 agent wallet and RPC billing;
    HIP-0095 QoS challenge;
    HIP-0096 AI compute contribution rewards;
    HIP-0106 plugin contract and the usage warehouse;
    HIP-1313 \code{/v1/usage};
    HIP-1325 \code{/v1/node}.
  \item[Lux proposals and papers.]
    LP-018 network security modes;
    LP-77 P-Chain warp payloads;
    LP-313 verifiable inference;
    LP-314 sovereign inference L1;
    LP-323 separating capacity, work and identity;
    LP-324 proof-carrying settlement for compute marketplaces;
    LP-325 linear budgets.
  \item[Hanzo papers.]
    \emph{Attested Compute in \code{hanzod}: Boot Measurement, Host-Asserted
    Evidence, and the Gap to Hardware Attestation}.
\end{description}

\section*{Copyright}

Copyright and related rights waived via CC0.

\end{document}
